\section{Il worm ILOVEYOU: Come funzionava?}

\subsection{Che tipo di programma era}

ILOVEYOU era uno \textbf{script} di circa 300 righe, scritto in \textbf{VBScript}. Uno script è un file di testo con una lista di istruzioni, che Windows legge ed esegue tramite un componente chiamato \textbf{Windows Script Host}, attivo di default su Windows 98, Me e 2000. Non serviva installare nulla: bastava un doppio clic sul file.

Si limitava a usare, in modo malevolo, strumenti normali del sistema: quelli per gestire i file, il registro di configurazione e il programma di posta Outlook. Per questo gli antivirus dell'epoca, abituati a riconoscere i programmi dannosi da ``impronte'' già note, lo individuarono in ritardo.

L'estensione \texttt{.vbs} del file era nascosta da un'impostazione di Windows, come descritto nella Sezione~\ref{subsec:estensione}.

\subsection{Struttura del codice}

Il codice era organizzato in piccole funzioni, ciascuna con un compito preciso: copiarsi, registrarsi all'avvio, spedire la posta, scorrere le unità del disco e alterare i file. Il programma principale le richiamava in sequenza. Questa struttura modulare rese il worm molto semplice da modificare: nei giorni successivi comparvero decine di varianti (Sezione~\ref{subsec:varianti}).

\subsection{Fase 1: copiarsi e restare nel sistema}\label{subsec:fase1}

Appena eseguito, lo script leggeva il proprio contenuto e ne salvava delle copie nelle cartelle di sistema di Windows, con nomi come \texttt{MSKernel32.vbs} e \texttt{Win32DLL.vbs}, simili a quelli dei file veri, perché l'utente non li eliminasse.

Poi scriveva alcune voci nel \textbf{registro di sistema}, il database delle impostazioni di Windows. Le voci erano nelle chiavi \texttt{Run} e \texttt{RunServices}, che contengono l'elenco dei programmi da avviare a ogni accensione. Così il worm ripartiva a ogni riavvio, anche dopo la cancellazione dell'allegato originale.

\subsection{Fase 2: diffondersi}\label{subsec:fase2}

Outlook permetteva ad altri programmi di usarlo tramite un'interfaccia chiamata \textbf{COM}, senza chiedere conferma all'utente. Lo script ne approfittava per:

\begin{enumerate}
    \item leggere tutti i contatti della rubrica;
    \item creare un messaggio per ciascuno, con oggetto ``ILOVEYOU'', un breve testo che invitava ad aprire la lettera e una copia del worm in allegato;
    \item inviarlo dall'account della vittima.
\end{enumerate}

Per ogni indirizzo raggiunto, il worm scriveva una voce nel registro, così da non spedire due volte lo stesso messaggio alla stessa persona. Questo controllo evitava di far scattare sospetti troppo presto, ma non limitava la diffusione complessiva.

Se sul computer c'era \textbf{mIRC}, un programma di chat molto diffuso allora, il worm creava anche una pagina web (\texttt{LOVE-LETTER-FOR-YOU.HTM}) contenente una copia del proprio codice, e modificava il file di configurazione \texttt{script.ini} del programma. Da quel momento chiunque entrasse in un canale con la vittima riceveva automaticamente il file.

\subsection{Fase 3: danneggiare i file}\label{subsec:fase3}

Lo script scorreva tutte le unità del computer, comprese quelle condivise in rete, entrando in ogni cartella e trattando i file in base all'estensione:

\begin{itemize}
    \item \textbf{script e pagine web} (\texttt{.vbs}, \texttt{.js}, \texttt{.css}, \texttt{.hta}): il contenuto veniva sostituito con il codice del worm, quindi i dati originali andavano persi;
    \item \textbf{immagini} (\texttt{.jpg}): sostituite e rinominate aggiungendo \texttt{.vbs};
    \item \textbf{musica} (\texttt{.mp3}): nascosta e rinominata, ma senza toccare i dati, quindi spesso recuperabile.
\end{itemize}

Sulle unità di rete il danno poteva estendersi ai file di un intero ufficio, perché bastava un solo computer infetto con accesso alle cartelle condivise.

\subsection{Fase 4: scaricare un programma per rubare le password}

Lo script cambiava la \textbf{pagina iniziale} di Internet Explorer, puntandola a un indirizzo che conteneva un file eseguibile, \texttt{WIN-BUGSFIX.exe}. Alla successiva apertura del browser il file veniva scaricato e impostato per l'avvio automatico. Era un \textbf{password stealer}, cioè un programma che raccoglie le password salvate sul computer e le invia all'autore.

Una volta scaricato il file, lo script reimpostava la pagina iniziale su una pagina vuota, per cancellare le tracce. Questa componente ebbe un impatto limitato, perché le pagine che distribuivano il file, ospitate dal provider filippino Sky Internet, furono rimosse in poco tempo.

\subsection{Perché funzionò}

Il worm riuscì a diffondersi perché in quegli anni coincidevano quattro condizioni:

\begin{itemize}
    \item Windows eseguiva gli script con un semplice doppio clic;
    \item Outlook lasciava che altri programmi leggessero la rubrica e inviassero posta senza avvisare;
    \item i server di posta non bloccavano gli allegati eseguibili;
    \item Windows 9x non separava i permessi dei programmi, quindi lo script poteva modificare liberamente file e registro.
\end{itemize}

Nessuna di queste era un difetto vero e proprio: erano funzioni pensate per comodità, ma troppo permissive.